\section{Graph operations, symmetry, and enumeration}\label{sec:ops}
\begin{proposition}[Products and line graphs]\leavevmode\label{prop:products}
Suppose $G,H$ are simple, with orders $n,h$, sizes $m,l$, and degrees $d_u,e_x$.
The following counts hold, with product vertices $(u,x)\in[n]\times[h]$.
\[
\begin{array}{lcc}
\toprule
\text{graph} & \text{degree at }(u,x)&\text{number of edges}\\
\midrule
G\square H &d_u+e_x&mh+ln\\
G\times H &d_u e_x&2ml\\
G\boxtimes H&d_u+e_x+d_u e_x&mh+ln+2ml\\
G[H]&h d_u+e_x&mh^2+nl\\
\bottomrule
\end{array}
\]
If both factors are nonempty and connected, Cartesian distance is the sum
of the factor distances. The line graph $L(G)$ has $m$ vertices and
$\sum_v\binom{d_v}{2}$ edges. Its vertex $uv$ has degree $d_u+d_v-2$.
\end{proposition}
\begin{proof}
In the Cartesian product change one coordinate along an edge, leaving the
other fixed. In the tensor product change both along edges independently.
Their disjoint union of neighbour lists is the strong product's neighbour
list. In the lexicographic product either change the first coordinate along
an edge, freely choosing the new second coordinate, or leave the first fixed
and move along an edge of $H$. These disjoint alternatives give the middle
column. Sum it over all $(u,x)$ and apply the handshaking identity to the two
factors for the right column. Any Cartesian walk projects to walks after
deleting stationary steps. Their lengths sum to the original length. Following
two shortest walks in succession realizes that lower bound. Two different
edges of a simple graph sharing a vertex share exactly one. Count pairs there,
or count the other edges at the two ends of a given $uv$.
\end{proof}

\begin{proposition}[Homomorphisms, subgraphs, and minors]\leavevmode\label{idx:operations-homomorphisms-subgraphs-and-minors}
A homomorphism is a table $f:[n]\to[h]$ sending edges to edges. Composition
preserves this condition, and $G$ has a $q$-colouring exactly when it admits
a homomorphism to $K_q$. A subgraph embedding is an injective edge-preserving
table, and is induced precisely when it also preserves nonedges at distinct
arguments. Isomorphism, containment as an ordinary or induced subgraph, and
minor containment are all decidable by PR procedures returning certificates.
\end{proposition}
\begin{proof}
Substitute the image of each edge into the next table. In $K_q$ two colours
are adjacent exactly when distinct. Scan all maps (at most $\max(2,h)^n$
codes), checking the desired row conditions. Isomorphism adds surjectivity
and reflection of adjacency, so its inverse is a bounded lookup. For a
minor of $G$ modelled on $H$, search for $h$ disjoint nonempty connected
branch sets in $[n]$, checking for a crossing edge for each edge of $H$.
Use codes below $2^{nh}$ and connectivity restricted to each mask. In an
actual sequence of deletions and contractions keep the sets of original
vertices represented by surviving vertices. Unions joined by an edge stay
connected, and every surviving edge has two original representatives in the
indicated sets. Conversely retain a spanning tree in each branch set and
one crossing edge for each required adjacency, delete everything else and
contract the retained trees. This constructs the minor and proves completeness
of the decision procedure.
\end{proof}

\begin{proposition}[Labelled, connected, and unlabelled graphs]\leavevmode\label{prop:burn}
Put $N=\binom n2$. There are $a_n=2^N$ labelled simple graphs on $[n]$,
$\binom Nm$ of size $m$, and $c_n$ connected graphs for $n\ge1$, where
\[
 c_n=a_n-\sum_{j=1}^{n-1}\binom{n-1}{j-1}c_j a_{n-j}.
\]
If $u_n$ counts isomorphism classes, and $o_2(\pi)$ counts cycles of $\pi$
on unordered vertex pairs, then
\[
 n!\,u_n=\sum_{\pi\in S_n} 2^{o_2(\pi)} .
\]
The numbers of loopless digraphs, oriented simple graphs and tournaments
on $[n]$ are respectively $4^N,3^N,2^N$.
\end{proposition}
\begin{proof}
The upper triangle has $N$ independent bits, and a mask of weight $m$ is
an $m$-subset. Pascal's recurrence counts it by splitting at the last slot.
For a nonempty graph choose the component of vertex zero. Its other labels
form a $(j-1)$-subset of $[n]\setminus\{0\}$, its induced graph is connected,
and the complement carries an arbitrary graph. Sorting either set supplies
a canonical labelling. The joining construction has no cross edges. The case
$j=n$ gives $c_n$, so the displayed subtraction is exact. For a fixed
permutation, an invariant graph chooses one bit for each orbit of pairs.
Count all $(\pi,G)$ with $\pi G=G$ by their graph or by their permutation.
Group the graph side using the least code in each orbit. For a representative
$R$, every orbit member $H$ has as many transporters $\pi R=H$ as stabilizers
of $R$: compose with the least transporter, which exists by finite search.
Its stabilizer has the same size as that of $H$, by conjugation. The fibres
over $H$ partition all $n!$ permutations, so each orbit contributes exactly
$n!$. Permutations are counted by successively selecting unused images.
Finally for each pair of labels the allowed directed states number four
(neither, either, both), three (exclude both) or two (exactly one).
\end{proof}

\begin{corollary}[Finite random graph counts and a Ramsey lower certificate]\leavevmode\label{idx:operations-finite-random-graph-counts-and-a-ramsey-lower-certifica}
In the uniform list of $2^N$ graphs the mean number of $k$-cliques for
$1\le k\le n$ is $\binom nk/2^{\binom k2}$. If
$2\binom nk<2^{\binom k2}$ there exists a graph on $n$ vertices with neither
a $k$-clique nor an independent $k$-set, and the least such mask is PR.
\end{corollary}
\begin{proof}
Fix the $k$ labels. Requiring every internal bit to equal one leaves
$2^{N-\binom k2}$ masks. Requiring zero leaves the same number. Count incidences
of a mask with a distinguished clique (or independent set) and sum over label
sets. Under the strict inequality the total number of incidences of either
kind is smaller than $2^N$. If every mask had at least one incidence the
sum would be at least $2^N$. A bounded scan therefore finds a zero count.
\end{proof}

\begin{theorem}[The Petersen graph]\leavevmode\label{thm:petersen}
Take the two-subsets of $[5]$ as vertices, with disjointness for adjacency.
This graph has order $10$, size $15$, diameter $2$, girth $5$, and strongly
regular parameters $(10,3,0,1)$. Its automorphism group is the full $S_5$
acting on the underlying labels, so it has order $120$. Its chromatic number
is three.
\end{theorem}
\begin{proof}
There are $\binom52$ vertices. For a fixed pair choose two of the remaining
three labels to find its three neighbours. Disjoint pairs use four labels
and have no common neighbour. Intersecting distinct pairs use three, leaving
exactly one common neighbour. These counts imply the diameter and exclude
triangles and four-cycles (opposite vertices of a four-cycle have two common
neighbours). With underlying labels $1,\ldots,5$, the sequence
$12,34,15,23,45,12$ is a five-cycle. An independent family consists of
pairwise intersecting pairs. If it contains $ab,ac$, every pair avoiding
$a$ must be $bc$, and no fourth distinct pair can meet all three of
$ab,ac,bc$. Thus the only independent families of size four are the five
full stars about a label. Automorphisms permute the stars, and the star
pair about $a,b$ intersects in the single vertex $ab$. Their action on
stars determines their action everywhere. Every permutation of the five
labels preserves disjointness, proving the group assertion. Colour the star
at label 1 first, the remaining pairs containing 2 second, and the three
pairs on $\{3,4,5\}$ third. Each class is intersecting. An odd cycle cannot
alternate two colours all the way round.
\end{proof}

\begin{proposition}[Intersection representation]\leavevmode\label{idx:operations-intersection-representation}
Every finite simple graph is an intersection graph of finite sets.
\end{proposition}
\begin{proof}
For the intersection representation, assign each vertex its set of incident
edge labels plus a private marker. Distinct vertices' sets intersect precisely
in their common edge if there is one. The markers distinguish isolated
vertices and leave all intersection tests between different vertices unchanged.
\end{proof}

\begin{theorem}[The graphic matroid and submodularity]\leavevmode\label{thm:matroid}
For a multigraph on $[n]$ the acyclic edge masks form a matroid: they contain
the empty mask, are closed under restriction, and if $I,J$ are acyclic with
$|I|<|J|$, one can extend $I$ by an edge of $J\setminus I$ while preserving
acyclicity. Its rank on $A$ is $r(A)=n-c(A)$, and
\[
 r(A\cup B)+r(A\cap B)\ \le\ r(A)+r(B).
\]
\end{theorem}
\begin{proof}
The first two assertions follow by the cycle test. If every edge of $J$ lies
inside the components of $I$, then in each such component, with $b$ vertices,
$J$ has at most $b-1$ edges because its restriction is a forest. Summing gives
$|J|\le n-c(I)=|I|$, contrary to the hypothesis. A crossing edge, found by
a scan, joins two components and creates no cycle. A maximal forest in $A$
spans exactly its components, giving the rank formula and a basis algorithm.
The increase in rank on inserting an edge equals one if its endpoints were
in different components and zero otherwise. For nested edge sets this
increment can only decrease on passage to the larger set. Start with
$A\cap B\subseteq B$ and insert, on both sides, the same successive edges
of $A\setminus B$. Adding the increment inequalities gives
$r(A)-r(A\cap B)\ge r(A\cup B)-r(B)$.
\end{proof}

\begin{theorem}[Reconstruction from the vertex-deleted deck]\leavevmode\label{thm:deck}
For a simple graph $G$ on $n\ge3$ vertices let its deck be the multiset of
unlabelled cards $G-v$. From the deck one can reconstruct the degree multiset,
counts of all smaller induced subgraph types, and, if $G$ is regular or
disconnected, the isomorphism type of $G$. Both special classes are recognizable.
\end{theorem}
\begin{proof}
Encode an unlabelled graph by the least adjacency mask among all its
relabellings. A card of order $n-1$ therefore has a canonically labelled
representative, with repetitions retained in the deck. For each graph $H$
of order $j<n$, write $s(H,G)$ for the count of $j$-vertex subsets inducing
$H$ up to isomorphism. A subset survives deletion in precisely $n-j$ cards,
so $(n-j)s(H,G)=\sum_v s(H,G-v)$. The same argument works for copies of
ordinary subgraphs on $j$ vertices. In particular the edge count $m$ is
the sum of the card edge counts divided by $n-2$, and $d(v)=m-|E(G-v)|$.
These divisions are exact by the counts. If the degrees all equal $d$,
choose any card. The neighbours of its missing vertex have degree $d-1$ there
and the other vertices degree $d$. Attach one vertex to the former group.
For $d=0$ simply add an isolated vertex.

A connected graph with at least two vertices has at least two connected
cards: removing either of two leaves of a spanning tree preserves the rest
of the tree. In a disconnected graph, a deletion can give a connected card
only if it removes an entire singleton component and exactly one component
then remains. For $n\ge3$ this can happen at most once. This recognizes
connectivity. In the disconnected case every component has order $<n$.
Consider all connected types $H$ of smaller order, in decreasing order of size,
and write $b_H$ for their unknown component multiplicities. A connected
induced copy of $H$ lies inside one component. Thus
\[
 b_H=s(H,G)-\sum_{\substack{C\ {\rm connected}\\|H|<|C|< n}}
                   b_C\,s(H,C).
\]
Among components of the same order an induced copy uses the whole component,
so only type $H$ contributes. Each difference is a nonnegative integer by
the counting identity. Their disjoint union reconstructs $G$. All types
can be listed by scanning masks up to $2^{\binom n2}$ and taking least codes.
This proves the stated special cases without a general reconstruction premise.
\end{proof}
